\documentclass[../main.tex]{subfiles}
\begin{document}

\begin{center}
    \Large{\textbf{TÓM TẮT NỘI DUNG}}\\
\end{center}
\vspace{1cm}
Trong thời đại công nghệ phát triển mạnh mẽ, dữ liệu (data) được xem như một nguồn tài nguyên vô cùng quan trọng và không thể thiếu. Tuy nhiên, theo thời gian, lượng dữ liệu phát sinh ngày càng lớn, đặc biệt là trong các hệ thống giám sát như camera an ninh. Điều này khiến các hệ quản trị cơ sở dữ liệu truyền thống dần bộc lộ nhiều hạn chế như khó mở rộng, hiệu suất thấp, và dễ bị mất dữ liệu khi xảy ra sự cố hệ thống.

Để khắc phục các vấn đề trên, hệ cơ sở dữ liệu phân tán ra đời, mang lại hiệu suất cao hơn, khả năng chịu lỗi tốt và dễ dàng mở rộng mà không làm gián đoạn hoạt động. Trong hệ thống phân tán, dữ liệu được lưu trữ trên nhiều máy tính khác nhau, có thể đặt ở cùng một vị trí hoặc phân tán tại nhiều nơi. Việc phân phối này không chỉ giúp đảm bảo tính nhất quán và toàn vẹn dữ liệu mà còn tăng khả năng phục hồi khi một hoặc nhiều máy chủ gặp sự cố.

Trong đồ án tốt nghiệp này, tôi xây dựng một hệ thống cơ sở dữ liệu phân tán dựa trên kiến trúc Master - Slave và triển khai trên hệ thống máy tính nhúng. Việc ứng dụng nền tảng nhúng giúp tối ưu chi phí, tiết kiệm không gian, đồng thời hướng đến khả năng triển khai thực tế trong các hệ thống camera an ninh tại các khu vực cần giám sát liên tục.

Đặc biệt, hệ thống không chỉ có khả năng lưu trữ dữ liệu video lớn từ camera mà còn được tích hợp thêm tính năng nhận diện hành động trong video. Tính năng này cho phép tự động phát hiện các hành vi bất thường hoặc nguy hiểm, từ đó nâng cao mức độ an toàn và tính thông minh cho hệ thống giám sát.

Mục tiêu của đồ án tốt nghiệp là xây dựng một hệ thống có khả năng hoạt động ổn định trên các thiết bị nhúng, đảm bảo tính nhất quán dữ liệu, dễ dàng mở rộng khi cần thiết và duy trì hiệu suất cao trong cả lưu trữ lẫn phân tích dữ liệu hình ảnh.
\end{document}